%! Trigonometry and Polar Coordinates — Unit 3, Lesson 3.13 — Homework (Answer Key)
\documentclass[10pt]{article}
\usepackage{precalculus-article}
\usepackage{precalculus-key}   % pulls in -boxes; do NOT also load -boxes
\newcommand{\TallMath}[1]{$\displaystyle #1\rule[-1.4em]{0pt}{3.2em}$}

\begin{document}

\pageheader{Unit 3, Lesson 3.13}{Homework --- Answer Key}
\namedateperiod

\vspace{0.10in}

\begin{notesbox}{Problems 1--7}

\textbf{1.}\quad Convert each polar coordinate pair to rectangular coordinates.

\vspace{4pt}
(a)\quad $(6,\, \pi/6)$:

\vspace{4pt}
$x = 6\cos(\pi/6) = $ \ans{$6 \cdot \frac{\sqrt{3}}{2} = 3\sqrt{3}$} \qquad
$y = 6\sin(\pi/6) = $ \ans{$6 \cdot \frac{1}{2} = 3$} \qquad
Rectangular: \ans{$(3\sqrt{3},\; 3)$}

\vspace{8pt}
(b)\quad $(4,\, 5\pi/4)$:

\vspace{4pt}
$x = $ \ans{$4\cos(5\pi/4) = 4 \cdot (-\frac{\sqrt{2}}{2}) = -2\sqrt{2}$} \qquad
$y = $ \ans{$4\sin(5\pi/4) = 4 \cdot (-\frac{\sqrt{2}}{2}) = -2\sqrt{2}$} \qquad
Rectangular: \ans{$(-2\sqrt{2},\; -2\sqrt{2})$}

\vspace{0.14in}
\textbf{2.}\quad Convert each rectangular point to polar form.

\vspace{4pt}
(a)\quad $(\sqrt{3},\, 1)$:

\vspace{4pt}
$r = $ \ans{$\sqrt{3+1} = 2$} \qquad
$\theta = $ \ans{$\arctan(1/\sqrt{3}) = \pi/6$} (quadrant: \ans{Q1}) \qquad
Polar: \ans{$(2,\; \pi/6)$}

\vspace{8pt}
(b)\quad $(-2,\, 2\sqrt{3})$:

\vspace{4pt}
$r = $ \ans{$\sqrt{4+12} = 4$} \qquad
$\theta = $ \ans{$\arctan(2\sqrt{3}/(-2)) + \pi = \arctan(-\sqrt{3}) + \pi = -\pi/3 + \pi = 2\pi/3$}
(quadrant: \ans{Q2}) \qquad
Polar: \ans{$(4,\; 2\pi/3)$}

\vspace{0.14in}
\textbf{3.}\quad $(0,-5)$: $r = 5$, $\theta = 3\pi/2$.

\vspace{4pt}
Rep.\ 1: \ans{$(5,\; 3\pi/2)$} \quad
Rep.\ 2: \ans{$(5,\; 3\pi/2 + 2\pi) = (5,\; 7\pi/2)$} \quad
Rep.\ 3: \ans{$(-5,\; \pi/2)$}

\vspace{0.14in}
\textbf{4.}\quad

\vspace{4pt}
(a)\quad $-3+3i$:

\vspace{4pt}
$r = $ \ans{$\sqrt{9+9} = 3\sqrt{2}$} \qquad
$\theta = $ \ans{$\arctan(3/(-3)) + \pi = -\pi/4 + \pi = 3\pi/4$} \qquad
Polar form: \ans{$3\sqrt{2}\cos(3\pi/4) + 3\sqrt{2}\,i\sin(3\pi/4)$}

\vspace{8pt}
(b)\quad $4i$:

\vspace{4pt}
$r = $ \ans{$4$} \qquad
$\theta = $ \ans{$\pi/2$} \qquad
Polar form: \ans{$4\cos(\pi/2) + 4i\sin(\pi/2)$}

\vspace{0.14in}
\textbf{5.}\quad $(2,\, \pi/4)$:

\vspace{4pt}
(a)\quad Rectangular: \ans{$(\sqrt{2},\; \sqrt{2})$}

\vspace{6pt}
(b)\quad Rep.\ 1: \ans{$(2,\; \pi/4 + 2\pi) = (2,\; 9\pi/4)$} \qquad
Rep.\ 2: \ans{$(-2,\; \pi/4 + \pi) = (-2,\; 5\pi/4)$}

\vspace{0.14in}
\textbf{6.}\quad \ans{(B)} \quad
$(-\sqrt{3},1)$: $r = \sqrt{3+1} = 2$; Q2 so $\theta = \arctan(-1/\sqrt{3})+\pi = -\pi/6+\pi = 5\pi/6$.

\vspace{0.12in}
\textbf{7.}\quad \ans{(C)} \quad
$(3,\, 2\pi/3)$: $-r$ with $\theta+\pi$ gives $(-3,\; 2\pi/3+\pi) = (-3,\; 5\pi/3)$.

\end{notesbox}

\vspace{0.08in}

\begin{tcolorbox}[colback=navylight, colframe=navy, arc=2mm,
    left=3mm, right=3mm, top=2mm, bottom=2mm]
\small\textbf{Preview --- Lesson 3.14: Polar Function Graphs}\\
You can now locate any point using $(r,\theta)$. In Lesson 3.14 we let $r$
be a \emph{function} of $\theta$ and graph equations like $r = 1+\cos\theta$.
The resulting curves --- cardioids, limaçons, rose curves --- are beautiful
and nearly impossible to express in rectangular form.
\end{tcolorbox}

\end{document}
